% Example document using qbook class — comprehensive guide
% Build from package root: make qbook (or make)
% SPDX-FileCopyrightText: 2026 Frank Langbein <frank@langbein.org>
% SPDX-License-Identifier: LPPL-1.3c
%
% This file demonstrates every qbook / qtex feature and serves as a
% compact user manual.  Lipsum text is used to fill pages so that
% headers, footers, spacing, and page layout are visible.
%
% Class options (try uncommenting alternatives):
%   Font:        sans (default) | serif | hacker
%   Layout:      oneside (default) | twoside
%   Citation:    numeric (default) | harvard
%   Alignment:   raggedright (default) | justified
%   Paragraph:   parskip (default) | parindent
%   Line spacing: single (default) | onehalfspacing | linespread115 | linespread12
%   Extras:      fullwidth, italic, obfuscate (LuaLaTeX)
%TC:subst \string\ STR

\documentclass[sans,oneside,numeric,raggedright,fullwidth]{qbook}
% \documentclass[serif,twoside,harvard,justified,parindent]{qbook}
% \documentclass[hacker,oneside,numeric,raggedright]{qbook}

\usepackage{lipsum}

% --------------------------------------------------------------------------
%  Metadata
% --------------------------------------------------------------------------
\title{The \texttt{qbook} Class --- User Guide}
\author{Author Name}
\university{Example Publishing House}
\date{\today}

\license{LPPL-1.3c}
\copyrightholder{Author Name}
\licenseyear{2026}
\keywords{example, book, qtex, typography, layout, guide}

% --------------------------------------------------------------------------
%  Document
% --------------------------------------------------------------------------
\begin{document}
\frontmatter

% \maketitle{abstract} prints: title page, copyright page (if \license is
% set), abstract chapter, and \frontlists (TOC, LOF, LOT, acronyms).
\maketitle{%
  This book serves as a comprehensive guide to the \texttt{qbook} class,
  which extends the standard \LaTeX{} \texttt{book} class with
  typography, layout, citations, and metadata features shared across the
  \texttt{qtex} package family.  It demonstrates every command and option
  the class provides and uses lipsum text so that page layout, headers,
  and spacing are visible across multiple pages.
}

% --------------------------------------------------------------------------
%  Dedication (unnumbered chapter in front matter)
% --------------------------------------------------------------------------
\dedication{To everyone who enjoys well-typeset books.}

% --------------------------------------------------------------------------
%  Preface
% --------------------------------------------------------------------------
\chapter{Preface}

This guide covers the structure and features of the \texttt{qbook} class.
Unlike the \texttt{qarticle} class (which is section-based), \texttt{qbook}
supports front matter, chapters, appendices, and back matter --- making it
suitable for theses, dissertations, monographs, and technical books.

The first chapter explains class options and document structure commands.
The second chapter demonstrates typography and layout features.  The third
chapter covers content elements such as figures, tables, lists,
mathematics, citations, and acronyms.  An appendix shows appendix
numbering.

% --------------------------------------------------------------------------
%  Acknowledgements
% --------------------------------------------------------------------------
\chapter{Acknowledgements}

The author thanks colleagues, friends, and institutions that supported the
creation of this guide.  Acknowledgements go here, in an unnumbered front
matter chapter.

\lipsum[1]

% ==========================================================================
\mainmatter
% ==========================================================================

% **************************************************************************
\chapter{Using \texttt{qbook}}\label{ch:using}
% **************************************************************************

\section{Class Options}\label{sec:options}

All options are forwarded to \texttt{qtex.sty}.  The \verb|\documentclass|
line at the top of this file uses the defaults; commented-out alternatives
show other combinations.

\subsection{Font families}

\begin{description}
  \item[\texttt{sans}] Source Sans Pro (default).
  \item[\texttt{serif}] Source Serif Pro.
  \item[\texttt{hacker}] Inconsolata monospaced.
\end{description}

\subsection{Layout and alignment}

\begin{description}
  \item[\texttt{oneside} / \texttt{twoside}] Single- or double-sided pages.
    With \texttt{twoside}, headers alternate and margins mirror.
  \item[\texttt{raggedright} / \texttt{justified}] Left-aligned (default) or
    fully justified text.
  \item[\texttt{parskip} / \texttt{parindent}] Paragraph separation by
    vertical space (default) or first-line indent.
  \item[\texttt{fullwidth}] Enables \texttt{fullwidth},
    \texttt{fullwidthfigure}, and \texttt{fullwidthtable} environments.
  \item[\texttt{italic}] Section headings and captions use italic instead of
    bold.
\end{description}

\subsection{Line spacing}

\begin{description}
  \item[\texttt{single}] Default.
  \item[\texttt{onehalfspacing}] 1.5 line spacing.
  \item[\texttt{linespread115}] \verb|\linespread{1.15}|.
  \item[\texttt{linespread12}] \verb|\linespread{1.2}|.
\end{description}

\subsection{Citation and other options}

\begin{description}
  \item[\texttt{numeric} / \texttt{harvard}] Numeric or author-year
    citations via \texttt{natbib}.
  \item[\texttt{obfuscate} / \texttt{noobfuscate}] Font obfuscation for
    blind review (LuaLaTeX only).
\end{description}

\lipsum[2]

\section{Document Structure}\label{sec:structure}

The \texttt{qbook} class follows the standard \LaTeX{} book structure
with some additions:

\subsection{Front matter}

\verb|\frontmatter| switches to roman page numerals.  Chapters in front
matter are unnumbered.  Do not use \verb|\subsection| in front matter (a
warning is issued).

\subsection{Title and abstract}

\verb|\maketitle{abstract}| prints the title page, an optional copyright
page (when \verb|\license{...}| is set), an abstract chapter with optional
keywords, and then calls \verb|\frontlists|.

\subsection{Front lists}

\verb|\frontlists| automatically generates the Table of Contents, and
conditionally the List of Figures, List of Tables, List of Algorithms, and
List of Acronyms --- only when each list is non-empty.

\subsection{Dedication}

\verb|\dedication{text}| prints a right-aligned, italicised dedication on
its own page.

\subsection{Main matter}

\verb|\mainmatter| switches to arabic page numerals and enables
chapter/section numbering (depth~3, i.e.\ down to \verb|\subsubsection|).

\subsection{Back matter and appendix}

\verb|\backmatter| restores appendix numbering (A, B, C, \ldots).
\verb|\appendix| marks the start of appendices.  Appendix chapters are
labelled ``Appendix A'', ``Appendix B'', etc., and appear in the table
of contents.

\subsection{Bibliography}

\verb|\bib{basename}| prints an optional License chapter (when
\verb|\license{...}| is set), then the bibliography.  The bibliography
style is chosen automatically:
\texttt{plainnat} for \texttt{numeric}, \texttt{agsm} for
\texttt{harvard}.

\lipsum[3]

\section{Metadata Commands}\label{sec:metadata}

\begin{itemize}
  \item \verb|\title{...}| --- book title.
  \item \verb|\author{...}| --- author name.
  \item \verb|\degree{...}| --- optional line under the author (e.g.\
    ``Second Edition'').
  \item \verb|\university{...}| --- affiliation or publisher.
  \item \verb|\date{...}| --- date.
  \item \verb|\keywords{...}| --- keywords shown in the abstract.
  \item \verb|\license{SPDX-ID}| --- enables copyright and license output.
  \item \verb|\copyrightholder{...}| --- copyright holder name(s).
  \item \verb|\licenseyear{YYYY}| --- copyright year.
\end{itemize}

% **************************************************************************
\chapter{Typography and Layout}\label{ch:typography}
% **************************************************************************

\section{Page Headers}

Running headers show the current chapter or section title in letterspaced
small caps, with the page number on the outer edge.  Chapter opening pages
use a special ``chapter'' page style with a thick rule and the chapter
label (``Chapter~1'', ``Appendix~A'').

\lipsum[4]

\section{Chapter and Section Headings}

Chapter titles are set in \verb|\Huge| bold (or italic with the
\texttt{italic} option).  Section headings use \verb|\Large| bold,
subsection \verb|\large| bold, and subsubsection \verb|\normalsize| bold.
Spacing above and below headings is configured by \texttt{titlesec} in
\texttt{qtex.sty}.

\subsection{A Subsection}

\lipsum[5]

\subsubsection{A Subsubsection}

\lipsum[6]

\section{New Thought}

\newthought{New thoughts} introduce a paragraph with letterspaced small
caps, creating visual rhythm inspired by Tufte's book designs.

\lipsum[7]

\section{Epigraph}

\epigraph{A book is a dream that you hold in your hand.}{Neil Gaiman}

An epigraph is a short quotation placed flush-right, typically at the
start of a chapter or section.

\section{Fullwidth Environments}

With the \texttt{fullwidth} class option, three environments extend into
the margins:

\begin{fullwidth}
  \textbf{This paragraph is inside a \texttt{fullwidth} environment.}
  It extends 2\,cm into both margins, giving a wider line for
  tables, figures, or text that benefits from extra horizontal space.
  \lipsum[8]
\end{fullwidth}

See also \texttt{fullwidthfigure} and \texttt{fullwidthtable} for floats
that automatically span the full width.

\section{Verse, Quotation, Quote}

\begin{verse}
  Shall I compare thee to a summer's day?\\
  Thou art more lovely and more temperate.
\end{verse}

\begin{quotation}
  Typography is a beautiful group of letters, not a group of beautiful
  letters.  A quotation block uses small text with a slight indent.
  \lipsum[9]
\end{quotation}

\begin{quote}
  Short quotes use the \texttt{quote} environment with no paragraph indent.
\end{quote}

% **************************************************************************
\chapter{Content Elements}\label{ch:content}
% **************************************************************************

\section{Lists}

\subsection{Itemize}

\begin{itemize}
  \item First item
    \begin{itemize}
      \item Sub-item one
      \item Sub-item two
        \begin{itemize}
          \item Sub-sub-item
        \end{itemize}
    \end{itemize}
  \item Second item
  \item Third item
\end{itemize}

\subsection{Enumerate}

\begin{enumerate}
  \item Step one
    \begin{enumerate}
      \item Sub-step A
      \item Sub-step B
    \end{enumerate}
  \item Step two
  \item Step three
\end{enumerate}

\subsection{Description}

\begin{description}
  \item[qbook] Book-style class with chapters and appendices.
  \item[qarticle] Article-style class with sections only.
  \item[qtex.sty] Shared typography, layout, and citation package.
\end{description}

\section{Mathematics}

Inline: $a^2 + b^2 = c^2$ and $\int_0^1 x\,\mathrm{d}x = \tfrac{1}{2}$.

Numbered:
\begin{equation}\label{eq:euler}
  e^{i\pi} + 1 = 0
\end{equation}

Aligned:
\begin{align}
  \nabla\cdot\mathbf{E} &= \frac{\rho}{\varepsilon_0}  \label{eq:gauss}\\
  \nabla\times\mathbf{B} &= \mu_0\mathbf{J} + \mu_0\varepsilon_0\frac{\partial\mathbf{E}}{\partial t} \label{eq:ampere}
\end{align}

See Equations~\eqref{eq:euler}--\eqref{eq:ampere}.

\section{Figures}

\begin{figure}[htbp]
  \centering
  \includegraphics[width=0.5\linewidth]{example-image-a}
  \caption{A placeholder figure.  Captions use \texttt{\string\small} and
    are configured by the \texttt{caption} package in \texttt{qtex.sty}.}
  \label{fig:example}
\end{figure}

Figure~\ref{fig:example} shows a placeholder image.

\begin{fullwidthfigure}[htbp]
  \includegraphics[width=\linewidth]{example-image-b}
  \caption{A fullwidth figure extending into both margins (requires the
    \texttt{fullwidth} class option).}
  \label{fig:fullwidth}
\end{fullwidthfigure}

\section{Tables}

\begin{table}[htbp]
  \centering
  \caption{Summary of runs.}\label{tab:summary}
  \begin{tabular}{lccc}
    \hline
    Run      & Mean & Std.\ dev. & $n$ \\
    \hline
    Baseline & 1.20 & 0.10 & 50 \\
    Var.\ A  & 1.52 & 0.21 & 50 \\
    Var.\ B  & 1.08 & 0.15 & 50 \\
    \hline
  \end{tabular}
\end{table}

Table~\ref{tab:summary} summarises the results.

\section{Footnotes}

Footnotes\footnote{This is a footnote.  The rule spans 40\% of the column
width.} are numbered per chapter and use compact spacing configured in
\texttt{qtex.sty}.

\section{Citations}

Use \verb|\citep{key}| for parenthetical: \citep{example-book}, and
\verb|\citet{key}| for textual: \citet{example-paper}.
Multiple keys: \citep{example-book,example-conf}.

\lipsum[10]

\section{Acronyms}

First use: \gls{api}.  Second use: \gls{api}.  Plural: \glspl{ui}.
The acronym list is printed automatically by \verb|\frontlists| when
acronyms exist.  Override the acronym file:

\noindent\verb|\renewcommand{\qtexacronymfile}{yourfile.tex}|

\lipsum[11]

% ==========================================================================
\backmatter
% ==========================================================================

\appendix
\chapter{Supplementary Material}\label{app:supplement}

Appendix chapters use letters: A, B, C, \ldots.  Sections within an
appendix are numbered A.1, A.2, etc.

\section{Extended Data}\label{app:data}

\lipsum[12]

\section{Derivations}\label{app:derivations}

\lipsum[13]

% --------------------------------------------------------------------------
%  Bibliography (with automatic License chapter when \license{...} is set)
% --------------------------------------------------------------------------
\bib{bibliography}
\end{document}
